\subsubsection{Alfabeto decimal y congruencia de control del emisor}
\label{sec:rec27-codigo}

Se conserva la ley emisora de \eqref{act21:tpk:bloque-decimal},
con las lecturas APP y el observable \(\ell_9\) de
\eqref{act21:tpk:logaritmo}, los cursores de
\ref{act21:tpk:operaciones} y el calendario de seis ventanas de
\ref{act21:tpk:emision}. Los acumuladores son \(S_m,E_m,Z_m\), con
\(Z_m=3\), y las cifras de salida se denotan \(d_{m,1},d_{m,2},d_{m,3}\).
La proposición~\ref{act21:tpk:recuperacion-firmas} recupera ya las dos
firmas residuales. Su combinación con las imágenes efectivas de los
acumuladores determina el alfabeto siguiente.

\begin{proposition}[Congruencia de control de la tercera cifra]
Cada bloque emitido satisface
\[
 d_{m,3}=[5d_{m,2}-2d_{m,1}+1]_{10}.
\]
\end{proposition}
\begin{proof}
La inversa de la proposición~\ref{act21:tpk:recuperacion-firmas}
da \([E_m]_{10}=[d_{m,2}-d_{m,1}]_{10}\). Sustituyendo
$E_m\equiv d_{m,2}-d_{m,1}\pmod{10}$ y $Z_m=3$ en la tercera cifra se obtiene
$3d_{m,1}+5(d_{m,2}-d_{m,1})+21\equiv5d_{m,2}-2d_{m,1}+1\pmod{10}$.
\end{proof}

\begin{proposition}[Alfabeto decimal de cuarenta y dos bloques]
La imagen del emisor por ventana consta exactamente de los bloques
\[
\begin{array}{c|rrrrrrr}
[E_m]_{10}&\multicolumn{7}{c}{U_m}\\ \hline
0&114&227&443&556&669&885&998\\
1&129&232&458&561&674&890&903\\
3&149&252&478&581&694&810&923\\
5&169&272&498&501&614&830&943\\
7&189&292&418&521&634&850&963\\
9&109&212&438&541&654&870&983
\end{array}
\]
\end{proposition}
\begin{proof}
Una trayectoria cardinal lee tres términos consecutivos del ciclo aditivo
\[
 (2,3,4,5,6,7,8,9,1).
\]
Sus sumas son
$(9,12,15,18,21,24,18,12,6)$; invertir la dirección conserva el conjunto
de sumas. Por tanto
\[
 \operatorname{Im}S_m=\{6,9,12,15,18,21,24\},\qquad
 \operatorname{Im}[S_m]_{10}=\{1,2,4,5,6,8,9\}.
\]
En el canal multiplicativo, un factor fijo divisible por tres aporta cero.
Para los factores invertibles, las sumas cíclicas de tres lecturas son
\[
\begin{array}{c|c}
\text{factor fijo}&\text{sumas posibles}\\\hline
1,8&\{1,3,7,9\}\\
2,7&\{3,5,9\}\\
4,5&\{1,5,7\}
\end{array}
\]
De aquí resulta $\operatorname{Im}E_m=\{0,1,3,5,7,9\}$.
La alcanzabilidad puede verificarse ya en la primera ventana: todos los
cursores de las dos tablas siguientes tienen dirección $E$.
\[
\begin{array}{c|rrrrrrr}
S_1&6&9&12&15&18&21&24\\\hline
(i_+,j_+)&(0,8)&(0,0)&(0,1)&(0,2)&(0,3)&(0,4)&(0,5)
\end{array}
\]
\[
\begin{array}{c|rrrrrr}
E_1&0&1&3&5&7&9\\\hline
(i_\times,j_\times)&(2,0)&(0,0)&(0,1)&(1,1)&(0,2)&(0,4)
\end{array}
\]
La independencia de ambos cursores realiza cada uno de los $7\cdot6$
pares. Las dos primeras cifras recuperan el par residual, de modo que sus
42 imágenes son distintas. La congruencia de control fija la tercera cifra
y produce exactamente la tabla enunciada.
\end{proof}

El bloque decimal conserva así las dos firmas residuales mediante una
codificación inyectiva y una cifra de control. Este resultado describe el
alfabeto del emisor finito; la prolongación de las lecturas regionales
conserva además las condiciones de fase, orientación y memoria del estado
enriquecido.

La imagen por ventana tiene, por tanto, \(42\) bloques. La concatenación
compatible de seis ventanas conserva el censo de \(468\) emisiones de
la proposición~\ref{act21:tpk:468}, cuya reducción ternaria produce
\(243\) palabras en el ambiente de \(729\). Cada cardinal cuenta el
objeto de su etapa y conserva la procedencia del mismo emisor.
